\documentclass[11pt]{article}
\usepackage{latexsym}
\oddsidemargin=-0.5cm
\evensidemargin=-0.5cm
\textwidth=16.80cm
\textheight=23.00cm
\topmargin=-2cm
\headsep=1cm

\title{(Agentic) Software Engineering\break BU CAS CS 413 Fall 2026}
\author{(Syllabus)}
\date{}

\begin{document}
\maketitle
\thispagestyle{empty}

\begin{itemize}
\item
{\bf Semester}
Fall, 2026
\item
{\bf Teaching Staff}
\begin{itemize}
\item
Instructor -- Hongwei Xi
\item
Teaching Assitant -- Victor Wang
\begin{itemize}
\item
Office Hours: TBA
\item
Location: CDS 1013, e-mail: \texttt{hwxi@cs.bu.edu}
\end{itemize}
\end{itemize}

\item {\bf Lecture Times}: TR 12:30-1:45PM
\item {\bf Classroom}: CAS 211
\item {\bf Reference Books}:\kern6pt
\begin{itemize}
%%%%%%
\item
Ian Sommerville, Software Engineering, {\bf 9th} edition,
Pearson/Addison-Wesley, 2011.
\begin{itemize}
\item ISBN: 978-0-13-703515-1
\item {\em The textbook provides the
conceptual foundation for the course. Some topics involving current
development practices will be supplemented with additional readings
and examples.}
\end{itemize}
%%%%%%
\item
Robert C. Martin, Clean Code: A Handbook of Agile Software Craftsmanship, Prentice Hall, 2008.
\begin{itemize}
\item 
ISBN: 978-0-13-235088-4
\item
{\em
This book supplements Sommerville by focusing more closely on the
design and implementation of maintainable source code.}
\end{itemize}
%%%%%%
\item
Erich Gamma, Richard Helm, Ralph Johnson, and John Vlissides, Design
Patterns: Elements of Reusable Object-Oriented Software,
Addison-Wesley, 1994.
\begin{itemize}
\item
ISBN: 978-0-201-63361-0
\item
{\em
This book presents a vocabulary and catalog of recurring solutions to
object-oriented software design problems.}
\end{itemize}
%%%%%%
\end{itemize}

\item {\bf Homepage}:
{\tt  https://hwxi.github.io/TEACHING/CS413/2026Fall}

\item {\bf Description}:
This course introduces the principles and practices used in the
engineering of medium/large-scale software systems. It focuses on
software development as a disciplined and collaborative engineering
activity rather than simply the construction of programs.

Primary topics include software processes, agile development,
requirements engineering, system modeling, architectural and
object-oriented design, software testing, software evolution,
dependability and security, software reuse, project management,
quality management, and configuration management.

A major component of the course is a semester-long team software
project (involving the construction of a compiler for a small
functional programming language). Students apply the techniques
discussed in class to the specification, design, implementation,
testing, and evolution of a nontrivial software system.

The course uses Ian Sommerville's Software Engineering, 9th edition,
as its primary text while supplementing the book with selected
contemporary software engineering practices.

\item
{\bf Course Philosophy}

Software engineering is fundamentally concerned with {\em complexity,
change, quality, and collaboration}.

Programming asks:
\begin{itemize}
\item How can we construct a program that performs a particular computation?
\end{itemize}

Software engineering asks a broader set of questions:
\begin{itemize}
\item What system should we build?
\item What do its users actually need?
\item How should those requirements be specified?
\item How should the system be structured?
\item How can multiple developers work on it safely?
\item How can we establish confidence that it works?
\item How do we manage changing requirements?
\item How do we manage defects and failures?
\item How do we release and maintain the system?
\item How can another developer understand and modify it years later?
\end{itemize}

The central goal of this course is therefore not merely to make
students better programmers. It is to teach students to approach
software development as an **engineering discipline** in which
requirements, design, evidence, tradeoffs, process, and communication
are as important as writing code.

\item {\bf Prerequisites}:
This course is designed for students who already have an immediate
level of proficiency in Python. You are expected to be familiar with a
text editor (e.g., vim, emacs, vscode).

\item {\bf Exams}:
\begin{itemize}
\item
First midterm exam: TBA
\item
Second midterm exam: TBA
\item
The final exam date is yet to be anounced.
\end{itemize}

\item {\bf Grades}
The final score is calculated using the following formula.
\[\mbox{final score = 30\%$\cdot$(homework) + 30\%$\cdot$(midterm) + 30\%$\cdot$(final) + 10\%$\cdot$(participation)}\]
The final letter grade is calculated as follows.
\begin{itemize}
\item{\bf A}: final score is $90\%$ or above
\item{\bf B}: final score is $80\%$ or above
\item{\bf C}: final score is $70\%$ or above
\item{\bf D}: final score is $60\%$ or above
\end{itemize}

\item
{\bf Homework Submission}:
Homework assignments are to be submitted via the Gradescope system.

\item
{\bf Attendance}:
It is expected that you will attend the lectures.

\item
{\bf Use of Generative AI}
Generative AI systems are increasingly used as software development
tools. Students may use such systems when explicitly permitted for an
assignment or project.

Permitted uses may include:
\begin{itemize}
\item Explaining unfamiliar code
\item Brainstorming alternative designs
\item Generating candidate test cases
\item Assisting with debugging
\item Generating routine implementation code
\item Reviewing code or documentation
\item Improving documentation
\end{itemize}

Students remain responsible for everything they submit. In particular,
students must be able to:
\begin{itemize}
\item Explain submitted code.
\item Explain important design decisions.
\item Verify generated code.
\item Test generated code.
\item Identify incorrect or inappropriate generated solutions.
\item Consider security and maintainability implications.
\item Properly acknowledge AI assistance when required.
\end{itemize}
The fact that code was generated by an AI system is not evidence that
the code is correct.

The instructor may ask a student to explain any submitted code or
design. Inability to explain one's own submission may affect the
student's grade and, where appropriate, raise an academic-integrity
concern.

\item
{\bf Academic Integrity}:
We adhere strictly to the standard BU guidelines for academic
integrity. For this course, it is perfectly acceptable for you to discuss
the general concepts and principles behind an assignment with other
students. However, it is not proper, without prior authorization of the
instructor, to arrive at collective solutions. In such a case, each student
is expected to develop, write up and hand in an individual solution and, in
doing so, gain a sufficient understanding of the problem so as to be able
to explain it adequately to the instructor.  Under {\em no} circumstances
should a student copy, partly or wholly, the completed solution of another
student. If one makes substantial use of certain code that is not written by
oneself, then the person must explicitly mention the source of the involved
code.

\end{itemize}

\end{document}
